% GENERATED FILE. Edit canonical lessons, then run npm run generate:books.
% canonical-source-hash: fnv1a64:869231b1310287da
% canonical-lessons: HI-C278-padhrahatha, HI-C278-pichla, HI-R278-the-scene-and-the-event

\chapter{What Was Going On}
\label{ch:hi-what-was-going-on}
\hlchaptermodality{278}

\begin{chapteropening}
\textbf{By the end of this chapter:} \emph{I can set a scene in the past with kar rahā thā and jab … tab, and say last week or last year with pichlā.}

Complete HI-R278-the-scene-and-the-event: set a scene in the past and say what happened inside it.
\end{chapteropening}

\section[\texorpdfstring{\hi{मैं} \hi{पढ़} \hi{रहा} \hi{था} (maiṁ paṛh rahā thā)}{मैं पढ़ रहा था (maiṁ paṛh rahā thā)}]{\hi{मैं} \hi{पढ़} \hi{रहा} \hi{था} (maiṁ paṛh rahā thā) — what was going on}

\label{lesson:HI-C278-padhrahatha}



\begin{quote}
Say \textquotedblleft{}I am doing\textquotedblright{} and \textquotedblleft{}he was\textquotedblright{}. (\emph{maiṁ kar rahā hūṁ}, \emph{vah thā}.) Then \textquotedblleft{}she came\textquotedblright{}. (\emph{vah āī}.) The first two make this lesson; the third will interrupt it.
\end{quote}

\subsection*{You'll want to know: was …-ing}
\begin{quote}
\textbf{\hi{मैं} \hi{किताब} \hi{पढ़} \hi{रहा} \hi{था।}}
\end{quote}

\begin{quote}
\emph{maiṁ kitāb paṛh rahā thā.} — \textquotedblleft{}I was reading a book.\textquotedblright{}
\end{quote}

You know \textbf{\hi{कर} \hi{रहा} \hi{हूँ}}, \textquotedblleft{}I am doing\textquotedblright{}. Put \textbf{\hi{था}} where \textbf{\hi{हूँ}} stood and the same action moves into the past. Both \textbf{\hi{रहा}} and \textbf{\hi{था}} follow the one doing it:

\begin{itemize}
\raggedright
  \item \textbf{\hi{मैं} \hi{पढ़} \hi{रहा} \hi{था}} (\emph{maiṁ paṛh rahā thā}) — I was reading; a woman says \textbf{\hi{मैं} \hi{पढ़} \hi{रही} \hi{थी}} (\emph{maiṁ paṛh rahī thī})
  \item \textbf{\hi{वह} \hi{सो} \hi{रहा} \hi{था}} (\emph{vah so rahā thā}) — he was sleeping; \textbf{\hi{वह} \hi{सो} \hi{रही} \hi{थी}} (\emph{vah so rahī thī}) — she was sleeping
  \item \textbf{\hi{बारिश} \hi{हो} \hi{रही} \hi{थी}} (\emph{bāriś ho rahī thī}) — it was raining
\end{itemize}

There is no \textbf{\hi{ने}} here, even with a thing: \textbf{\hi{मैं} \hi{खाना} \hi{खा} \hi{रहा} \hi{था}}, never \textbf{\hi{मैंने}}. \textbf{\hi{ने}} belongs to the plain past alone.

\begin{grammarlens}[title={Grammar Lens: a scene and an event}]
\begin{quote}
\textbf{\hi{जब} \hi{मैं} \hi{खाना} \hi{खा} \hi{रहा} \hi{था}, \hi{तब} \hi{मीरा} \hi{आई।}}
\end{quote}

\begin{quote}
\emph{jab maiṁ khānā khā rahā thā, tab Mīrā āī.} — \textquotedblleft{}When I was eating, Mira came.\textquotedblright{}
\end{quote}

A story needs two kinds of past. \textbf{\hi{खा} \hi{रहा} \hi{था}} is the scene: it was going on, and the sentence does not say when it began or ended. \textbf{\hi{आई}} is the event: it happened, once, inside the scene. \textbf{\hi{जब} … \hi{तब} …}, the frame you know for \textquotedblleft{}when … then …\textquotedblright{}, holds the two together.

\begin{itemize}
\raggedright
  \item \textbf{\hi{जब} \hi{बारिश} \hi{हो} \hi{रही} \hi{थी}, \hi{तब} \hi{बिजली} \hi{गई।}} (\emph{jab bāriś ho rahī thī, tab bijlī gaī.}) — While it was raining, the power went.
\end{itemize}
\end{grammarlens}

\subsection*{Your turn}
Say these aloud:

\begin{itemize}
\raggedright
  \item \textquotedblleft{}What were you doing yesterday evening?\textquotedblright{} — \emph{tum kal śām ko kyā kar rahe the?}
  \item \textquotedblleft{}We were drinking tea when the power went\textquotedblright{} — \emph{jab ham cāy pī rahe the, tab bijlī gaī}
\end{itemize}

\begin{tcolorbox}[breakable,colback=teal!4,colframe=teal!35!black,title={Before you move on}]
\textquotedblleft{}I was reading\textquotedblright{}? (\textbf{\hi{मैं} \hi{पढ़} \hi{रहा} \hi{था}}, or \textbf{\hi{पढ़} \hi{रही} \hi{थी}}.) \textquotedblleft{}It was raining\textquotedblright{}? (\textbf{\hi{बारिश} \hi{हो} \hi{रही} \hi{थी}}.) In \textbf{\hi{जब} \hi{मैं} \hi{खाना} \hi{खा} \hi{रहा} \hi{था}, \hi{तब} \hi{मीरा} \hi{आई}}, which is the scene and which is the event? (\textbf{\hi{खा} \hi{रहा} \hi{था}} is the scene; \textbf{\hi{आई}} is the event.) Why no \textbf{\hi{ने}} in \textbf{\hi{मैं} \hi{खाना} \hi{खा} \hi{रहा} \hi{था}}? (\textbf{\hi{ने}} belongs to the plain past alone.)
\end{tcolorbox}
\section[\texorpdfstring{\hi{पिछला} (pichlā)}{पिछला (pichlā)}]{\hi{पिछला} (pichlā) — last, previous}

\label{lesson:HI-C278-pichla}



\begin{quote}
Before the new one: say \textquotedblleft{}I was reading a book\textquotedblright{}. (\emph{maiṁ kitāb paṛh rahā thā}.) Then the words for a week and a year. (\emph{haftā}, \emph{sāl}.)
\end{quote}

\subsection*{You'll want to know: \hi{पिछला}}
\textbf{\hi{पिछला}} — \emph{pichlā} — \textquotedblleft{}last, previous\textquotedblright{}: the one before this one.

\begin{quote}
\textbf{\hi{पिछले} \hi{हफ़्ते} \hi{मैं} \hi{गाँव} \hi{गया।}}
\end{quote}

\begin{quote}
\emph{pichle hafte maiṁ gā̃v gayā.} — \textquotedblleft{}Last week I went to the village.\textquotedblright{}
\end{quote}

\textbf{\hi{पिछला}} is an adjective, and it agrees: \textbf{\hi{पिछला} \hi{हफ़्ता}} (\emph{pichlā haftā}), \textquotedblleft{}last week\textquotedblright{}, but \textbf{\hi{पिछली} \hi{बार}} (\emph{pichlī bār}), \textquotedblleft{}last time\textquotedblright{}, because \textbf{\hi{बार}} is feminine.

Said as a time, \textquotedblleft{}last week\textquotedblright{} or \textquotedblleft{}last year\textquotedblright{} takes the form you met in \textbf{\hi{बच्चे} \hi{को}}: \textbf{\hi{पिछले} \hi{हफ़्ते}}, \textbf{\hi{पिछले} \hi{साल}}.

\begin{quote}
\textbf{\hi{पिछले} \hi{साल} \hi{हम} \hi{बहुत} \hi{व्यस्त} \hi{थे।}}
\end{quote}

\begin{quote}
\emph{pichle sāl ham bahut vyast the.} — \textquotedblleft{}Last year we were very busy.\textquotedblright{}
\end{quote}

\begin{cousinweb}
\textbf{\hi{पिछला}} (\emph{pichlā}) has no row of its own in CDIAL. It is usually connected with \textbf{\hi{पीछे}} (\emph{pīche}), \textquotedblleft{}behind\textquotedblright{}, which CDIAL traces to \emph{*paśca}, \textquotedblleft{}hinder part\textquotedblright{} (CDIAL 7990). The last week, on that picture, is the one behind you.
\end{cousinweb}

\subsection*{Your turn}
Say these aloud:

\begin{itemize}
\raggedright
  \item \emph{pichlā}, then \emph{pichle hafte}
  \item \textquotedblleft{}Last week I went to the village\textquotedblright{} — \emph{pichle hafte maiṁ gā̃v gayā}
  \item \textquotedblleft{}Last time it was raining\textquotedblright{} — \emph{pichlī bār bāriś ho rahī thī}
\end{itemize}

\begin{tcolorbox}[breakable,colback=teal!4,colframe=teal!35!black,title={Before you move on}]
What does \textbf{\hi{पिछला}} mean? (\textquotedblleft{}Last, previous\textquotedblright{}.) \textquotedblleft{}Last year\textquotedblright{}? (\textbf{\hi{पिछले} \hi{साल}}.) \textquotedblleft{}Last time\textquotedblright{}? (\textbf{\hi{पिछली} \hi{बार}}: \textbf{\hi{बार}} is feminine.)
\end{tcolorbox}
\section[(review)]{Review — the scene and the event}

\label{lesson:HI-R278-the-scene-and-the-event}



\begin{quote}
Say \textquotedblleft{}I was reading a book\textquotedblright{}. (\emph{maiṁ kitāb paṛh rahā thā}.) Then \textquotedblleft{}last week\textquotedblright{} and \textquotedblleft{}last time\textquotedblright{}. (\emph{pichle hafte}, \emph{pichlī bār}.)
\end{quote}

\subsection*{Your turn}
Say these aloud:

\begin{itemize}
\raggedright
  \item Last week I went to the village.
  \item What were you doing yesterday evening? I was cooking.
  \item When we were drinking tea, Arun came.
  \item It was raining, and suddenly the power went.
  \item Last year Mira read many books.
\end{itemize}

\begin{tcolorbox}[breakable,colback=teal!4,colframe=teal!35!black,title={Before you move on}]
Which past is the scene? (\textbf{\hi{रहा} \hi{था}}: \textbf{\hi{बारिश} \hi{हो} \hi{रही} \hi{थी}}.) Which is the event? (The plain past: \textbf{\hi{मीरा} \hi{आई}}.) \textquotedblleft{}Last year she read many books\textquotedblright{}? (\textbf{\hi{पिछले} \hi{साल} \hi{उसने} \hi{बहुत} \hi{किताबें} \hi{पढ़ीं}}: \textbf{\hi{ने}}, and \textbf{\hi{पढ़ीं}} for the books.)
\end{tcolorbox}
